% Referee Report -- Panel Adjudication Round 1
% Generated by gpd-referee on 2026-03-18
% Aligned with .gpd/REFEREE-REPORT.md

\documentclass[11pt]{article}

\usepackage[margin=1in]{geometry}
\usepackage[T1]{fontenc}
\usepackage[utf8]{inputenc}
\usepackage{lmodern}
\usepackage{microtype}
\usepackage[dvipsnames]{xcolor}
\usepackage{hyperref}
\usepackage{booktabs}
\usepackage{longtable}
\usepackage{tabularx}
\usepackage{array}
\usepackage{enumitem}
\usepackage{titlesec}
\usepackage{fancyhdr}
\usepackage[most]{tcolorbox}

\definecolor{ReportInk}{HTML}{12263A}
\definecolor{ReportCrimson}{HTML}{9B2226}
\definecolor{ReportGold}{HTML}{C76D1D}
\definecolor{ReportSage}{HTML}{5E7D64}
\definecolor{ReportStone}{HTML}{F6F1EA}
\definecolor{ReportLine}{HTML}{D9D2C8}

\hypersetup{
  colorlinks=true,
  linkcolor=ReportCrimson,
  urlcolor=ReportCrimson,
  citecolor=ReportCrimson,
  pdftitle={Referee Report: When do Vision Transformers help?},
  pdfauthor={gpd-referee}
}

\setlength{\parindent}{0pt}
\setlength{\parskip}{0.55em}
\setlist[itemize]{leftmargin=1.4em, itemsep=0.35em, topsep=0.35em}
\setlist[enumerate]{leftmargin=1.6em, itemsep=0.35em, topsep=0.35em}
\renewcommand{\arraystretch}{1.2}

\titleformat{\section}{\Large\bfseries\color{ReportInk}}{\thesection}{0.6em}{}
\titleformat{\subsection}{\large\bfseries\color{ReportInk}}{\thesubsection}{0.6em}{}
\titlespacing*{\section}{0pt}{1.4em}{0.6em}
\titlespacing*{\subsection}{0pt}{1.0em}{0.4em}

\pagestyle{fancy}
\fancyhf{}
\fancyhead[L]{\textsc{Mock Referee Report}}
\fancyhead[R]{\textcolor{ReportCrimson}{ViT vs CNN Glitch Classification}}
\fancyfoot[C]{\thepage}
\setlength{\headheight}{14pt}

\newcolumntype{Y}{>{\raggedright\arraybackslash}X}

\newcommand{\RecommendationBadge}[2]{%
  \tcbox[
    on line,
    box align=base,
    colback=#1!12,
    colframe=#1,
    arc=2pt,
    boxrule=0.8pt,
    left=7pt,
    right=7pt,
    top=4pt,
    bottom=4pt
  ]{\textbf{\textsc{#2}}}%
}

\newtcolorbox{SummaryBox}{
  enhanced,
  breakable,
  colback=ReportStone,
  colframe=ReportInk,
  boxrule=0.8pt,
  arc=2pt,
  left=12pt,
  right=12pt,
  top=10pt,
  bottom=10pt
}

\newtcolorbox{StrengthBox}{
  enhanced,
  breakable,
  colback=ReportSage!10,
  colframe=ReportSage,
  colbacktitle=ReportSage,
  coltitle=white,
  title={Strengths},
  boxrule=0.8pt,
  arc=2pt,
  left=12pt,
  right=12pt,
  top=10pt,
  bottom=10pt
}

\newtcolorbox{MajorIssueBox}[1]{
  enhanced,
  breakable,
  colback=white,
  colframe=ReportCrimson,
  colbacktitle=ReportCrimson,
  coltitle=white,
  title={#1},
  boxrule=0.9pt,
  arc=2pt,
  left=12pt,
  right=12pt,
  top=10pt,
  bottom=10pt
}

\newtcolorbox{MinorIssueBox}[1]{
  enhanced,
  breakable,
  colback=white,
  colframe=ReportGold,
  colbacktitle=ReportGold,
  coltitle=white,
  title={#1},
  boxrule=0.9pt,
  arc=2pt,
  left=12pt,
  right=12pt,
  top=10pt,
  bottom=10pt
}

\begin{document}

{\color{ReportCrimson}\rule{\linewidth}{1.1pt}}

\begin{center}
  {\Huge\bfseries\color{ReportInk} Referee Report\par}
  \vspace{0.45em}
  {\Large When do Vision Transformers help? Class-dependent architecture\\preferences for gravitational-wave glitch classification\par}
  \vspace{0.8em}
  \RecommendationBadge{ReportGold}{Major Revision}
\end{center}

\vspace{0.8em}

\begin{tcolorbox}[
  enhanced,
  colback=white,
  colframe=ReportLine,
  boxrule=0.7pt,
  arc=2pt,
  left=12pt,
  right=12pt,
  top=10pt,
  bottom=10pt
]
\begin{tabularx}{\linewidth}{@{}p{0.23\linewidth}Y p{0.23\linewidth}Y@{}}
\textbf{Scope} & Manuscript & \textbf{Target journal} & Classical and Quantum Gravity \\
\textbf{Reviewed} & 2026-03-18 & \textbf{Confidence} & high \\
\textbf{Major issues} & 5 & \textbf{Minor issues} & 8 \\
\textbf{Round} & Initial review & \textbf{Reviewer} & gpd-referee (panel adjudicator) \\
\end{tabularx}
\end{tcolorbox}

% ============================================================
% SUMMARY
% ============================================================

\begin{SummaryBox}
\textbf{Summary}

This paper presents a controlled comparison of ViT-B/16 and ResNet-50v2 BiT for 23-class gravitational-wave glitch classification on LIGO O3 Gravity Spy data. ViT significantly improves overall macro-F1 (0.723 vs.\ 0.679, $p = 0.0002$), but the rare-class comparison is statistically underpowered (aggregate power $= 0.20$). A secondary contribution connects classifier performance to CW search data quality through a sample-count-based duty-cycle proxy. Both models are evaluated zero-shot on O4 data, where the per-class pattern does not replicate as a monotonic sample-size trend.

The paper has genuine merit as a CQG detector-characterization methods paper. The identical-recipe comparison, honest negative results, temporal-split methodology, and CW veto connection are all useful contributions. However, the manuscript currently overclaims in three areas: (1) the rare-class CNN advantage is presented as directional in the introduction and conclusion despite being underpowered; (2) physical-interpretation claims about self-attention and inductive bias lack direct evidence and are asymmetrically hedged; (3) the novelty framing needs narrowing against Wu et al.\ (CQG 2025). All issues are fixable through targeted text revision without new experiments.

The mathematics and numerics are sound: all key quantities verified from source JSON to machine precision, with one minor double-rounding error (ViT O4 macro-F1: 0.670 should be 0.669).
\end{SummaryBox}

% ============================================================
% STRENGTHS
% ============================================================

\begin{StrengthBox}
\begin{enumerate}
  \item \textbf{Statistical hygiene above field standard.} Power analysis for rare-class comparison, paired bootstrap with 10K resamples, and temporal-split validation with random-split control are practices most GW ML papers omit.
  \item \textbf{Honest negative results.} O4 non-replication of per-class patterns and non-significant Spearman correlation reported straightforwardly.
  \item \textbf{CW veto connection.} Duty-cycle analysis connects glitch classification to downstream CW search sensitivity -- bridging a gap between classification and search-impact papers.
  \item \textbf{Temporal split methodology.} GPS-time split with 60-second gaps and random-split control quantify data-leakage effects for the first time in the Gravity Spy literature.
  \item \textbf{Reproducibility infrastructure.} All 37 number macros from single JSON source, full execution pipeline with seeds and checksums.
  \item \textbf{Number pipeline integrity.} All key quantities verified to machine precision by the math review stage.
\end{enumerate}
\end{StrengthBox}

% ============================================================
% MAJOR ISSUES
% ============================================================

\section*{Major Issues}

\begin{MajorIssueBox}{REF-001: Rare-class narrative overclaimed in Introduction and Conclusion}
\textbf{Dimension:} significance / correctness

\textbf{Location:} paper/main.tex:193, paper/main.tex:693--695

\textbf{Description:} The introduction states ``yet rare-class performance declines'' without qualification. The conclusion states ``CNN retains advantages on rare classes where limited training data appears insufficient for learning effective attention patterns,'' converting an underpowered observation ($p = 0.884$, power $= 0.20$) into a causal narrative. The abstract is appropriately hedged; the introduction and conclusion are not.

\textbf{Impact:} The headline ``class-dependent architecture preference'' claim partially depends on rare classes favoring CNN. If the rare-class direction is uncertain, the story must be honestly reframed throughout.

\textbf{Suggested fix:} Harmonize all rare-class language with the abstract: ``insufficient evidence to determine whether ViT improves or degrades rare-class performance.'' Condition the inductive-bias regularization interpretation on future confirmation.
\end{MajorIssueBox}

\begin{MajorIssueBox}{REF-002: Novelty claim too broad against Wu et al.\ (CQG 2025)}
\textbf{Dimension:} novelty / literature\_context

\textbf{Location:} paper/main.tex:167--170

\textbf{Description:} The manuscript states ``none of these studies has conducted a controlled per-class comparison under identical training conditions.'' Wu et al.\ (CQG 42, 165015) report class-wise accuracy tables and confusion matrices. The present paper's novelty is the identical-recipe paired comparison with statistical testing and power analysis, not the per-class analysis itself.

\textbf{Impact:} A CQG referee familiar with Wu et al.\ will flag this overclaim.

\textbf{Suggested fix:} Narrow to: ``none of these studies has conducted an identical-recipe paired comparison with statistical testing of per-class differences.'' Acknowledge Wu et al.\ per-class results.
\end{MajorIssueBox}

\begin{MajorIssueBox}{REF-003: Missing canonical review citation (Cuoco et al.\ 2025)}
\textbf{Dimension:} literature\_context

\textbf{Location:} paper/main.tex:130--170

\textbf{Description:} Cuoco, Cavaglia et al., \emph{Living Reviews in Relativity} 28, 2 (2025), is the current canonical review of ML for GW research. It is not cited. Wu et al.\ is cited as arXiv but now published in CQG.

\textbf{Impact:} Missing a canonical review in the target journal's subject area signals inadequate literature engagement.

\textbf{Suggested fix:} Add Cuoco et al.\ (2025) to the introduction. Update Wu et al.\ to CQG 42, 165015.
\end{MajorIssueBox}

\begin{MajorIssueBox}{REF-004: Physical-interpretation claims asymmetrically hedged}
\textbf{Dimension:} technical\_soundness / completeness

\textbf{Location:} paper/main.tex:339--340, paper/main.tex:523--534

\textbf{Description:} The Power\_Line self-attention interpretation is presented as explanatory while the rare-class CNN interpretation is hedged as ``speculative.'' Neither has direct evidence (no attention maps, ablations, or controlled experiments).

\textbf{Impact:} Asymmetric hedging creates an impression of established mechanism where only pattern-matching exists.

\textbf{Suggested fix:} Apply consistent ``speculative/hypothetical'' qualification to all inductive-bias interpretations. Alternatively, add attention-map visualizations.
\end{MajorIssueBox}

\begin{MajorIssueBox}{REF-005: CW abstract conflates veto efficiency with duty cycle}
\textbf{Dimension:} technical\_soundness / significance

\textbf{Location:} paper/main.tex:117--118, paper/main.tex:480--483

\textbf{Description:} The abstract reports ``approximately equal'' veto efficiency without mentioning the statistically significant CNN duty-cycle advantage ($\Delta_\mathrm{DC} = -0.051$, CI [$-0.054$, $-0.048$]). ``First quantitative assessment'' overstates a sample-count proxy.

\textbf{Impact:} CW practitioners would conclude architectures are interchangeable, missing the 5\,pp duty-cycle difference.

\textbf{Suggested fix:} Abstract must mention duty-cycle difference. Replace ``first quantitative assessment'' with ``first proxy-level comparison.''
\end{MajorIssueBox}

% ============================================================
% MINOR ISSUES
% ============================================================

\section*{Minor Issues}

\begin{MinorIssueBox}{REF-006: Double-rounding error in ViT O4 macro-F1}
\textbf{Dimension:} correctness

\textbf{Location:} paper/main.tex:55, table\_overall.tex:11

\textbf{Description:} Source value 0.6694683 rounds to 0.669, not the reported 0.670. Classic double-rounding.

\textbf{Suggested fix:} Correct \texttt{\textbackslash vitMacroFOfour} from 0.670 to 0.669.
\end{MinorIssueBox}

\begin{MinorIssueBox}{REF-007: ``8 rare classes'' labeling inconsistency}
\textbf{Dimension:} clarity

\textbf{Location:} paper/main.tex:354

\textbf{Description:} Paper defines rare as $N_\mathrm{train} < 200$ (4 classes) but power analysis tests 8 classes including non-rare ones.

\textbf{Suggested fix:} Change ``all 8 rare classes'' to ``all 8 classes with small test sets.''
\end{MinorIssueBox}

\begin{MinorIssueBox}{REF-008: Section heading overstates absence of correlation}
\textbf{Dimension:} correctness / clarity

\textbf{Location:} paper/main.tex:389

\textbf{Description:} Spearman test with $n=23$ has power $\approx 0.08$ to detect $\rho = -0.12$. Absence of evidence $\neq$ evidence of absence.

\textbf{Suggested fix:} Soften to ``No significant evidence that architecture preference is a monotonic function of sample size.''
\end{MinorIssueBox}

\begin{MinorIssueBox}{REF-009: 20\% O4 degradation threshold unmotivated}
\textbf{Dimension:} technical\_soundness

\textbf{Location:} paper/main.tex:448--450

\textbf{Description:} ViT degrades 7.4\% vs.\ CNN 1.7\% (4.4$\times$ ratio) but both ``pass'' an ad-hoc 20\% bar.

\textbf{Suggested fix:} Motivate from literature or remove pass/fail framing.
\end{MinorIssueBox}

\begin{MinorIssueBox}{REF-010: Forbidden-proxy lesson not cited from imbalanced-learning literature}
\textbf{Dimension:} literature\_context / novelty

\textbf{Location:} paper/main.tex:592--607

\textbf{Description:} Per-class evaluation under class imbalance is established (He \& Garcia 2009; Buda et al.\ 2018).

\textbf{Suggested fix:} Add 1--2 imbalanced-learning citations. Reframe as GW demonstration of known principle.
\end{MinorIssueBox}

\begin{MinorIssueBox}{REF-011: ViT interpretability literature not cited}
\textbf{Dimension:} literature\_context

\textbf{Location:} paper/main.tex:523--534

\textbf{Description:} Attention-mechanism claims not grounded in ViT interpretability literature.

\textbf{Suggested fix:} Cite Raghu et al.\ (2021) or similar.
\end{MinorIssueBox}

\begin{MinorIssueBox}{REF-012: One-sided $p$-value with opposite-direction effect}
\textbf{Dimension:} correctness

\textbf{Location:} table\_overall.tex:10

\textbf{Description:} One-sided ViT $>$ CNN gives $p = 0.884$ when ViT $<$ CNN. Two-sided $p = 0.232$, still non-significant.

\textbf{Suggested fix:} Switch to two-sided test or add footnote.
\end{MinorIssueBox}

\begin{MinorIssueBox}{REF-013: Wu et al.\ cited as arXiv, now published in CQG}
\textbf{Dimension:} presentation\_quality

\textbf{Location:} paper/main.tex:774--777

\textbf{Description:} Wu et al.\ is now CQG 42, 165015.

\textbf{Suggested fix:} Update citation.
\end{MinorIssueBox}

% ============================================================
% SUGGESTIONS
% ============================================================

\section*{Suggestions}
\begin{itemize}
  \item \textbf{Multi-seed training (3--5 seeds).} Would substantially strengthen per-class claims, especially for rare classes. Estimated effort: medium.
  \item \textbf{Attention-map visualizations.} Even without ablation, showing what ViT attends to on Power\_Line spectrograms would provide direct evidence. Estimated effort: small.
  \item \textbf{Restructure narrative for CQG.} Foreground CW veto findings; the F1-vs-DC disconnect is more CQG-relevant. Estimated effort: small.
  \item \textbf{Two-sided bootstrap tests.} Avoids awkward $p$ near 1.0. Estimated effort: trivial.
\end{itemize}

% ============================================================
% DETAILED EVALUATION
% ============================================================

\section*{Detailed Evaluation}

\begin{longtable}{@{}p{0.18\linewidth}p{0.17\linewidth}Y@{}}
\toprule
\textbf{Dimension} & \textbf{Rating} & \textbf{Assessment} \\
\midrule
\endhead
Novelty & ADEQUATE & Identical-recipe comparison with statistical testing is genuine; needs narrowing against Wu et al. \\
Correctness & MOSTLY CORRECT & All key quantities verified; one minor double-rounding error (REF-006). \\
Clarity & GOOD & Well-organized; inconsistent hedging between abstract and intro/conclusion. \\
Completeness & MOSTLY COMPLETE & All promised results delivered; single-seed limitation not propagated into claims. \\
Significance & ADEQUATE & Adequate for CQG as methods paper after reframing; would not survive more selective venue. \\
Reproducibility & FULLY REPRODUCIBLE & Seeds, checksums, full pipeline documented. \\
Literature Context & INCOMPLETE & Missing Cuoco et al.\ 2025; Wu et al.\ outdated; imbalanced-learning lit not cited. \\
Presentation Quality & NEEDS POLISHING & Good structure; claim inconsistencies and minor numerical error. \\
Technical Soundness & MOSTLY SOUND & Appropriate methodology; single seed is main concern. \\
Publishability & Major revision & Publishable at CQG after targeted claim recalibration. \\
\bottomrule
\end{longtable}

% ============================================================
% PHYSICS CHECKLIST
% ============================================================

\section*{Physics Checklist}

\begin{tabularx}{\linewidth}{@{}p{0.38\linewidth}p{0.18\linewidth}Y@{}}
\toprule
\textbf{Check} & \textbf{Status} & \textbf{Notes} \\
\midrule
Dimensional analysis & pass & All metrics dimensionless ratios. \\
Limiting cases & pass & Random-split control recovers published accuracy range. \\
Symmetry preservation & N/A & ML benchmark. \\
Conservation laws & N/A & ML benchmark. \\
Error bars present & pass & 95\% bootstrap CIs on all primary metrics. \\
Approximations justified & pass & Stationarity assumption stated; proxy acknowledged as coarse. \\
Convergence demonstrated & partial & Bootstrap 10K adequate; no multi-seed convergence. \\
Literature comparison & partial & CNN accuracy gap explained; Wu et al.\ per-class comparison missing. \\
Reproducible & pass & Full pipeline with seeds and checksums. \\
\bottomrule
\end{tabularx}

% ============================================================
% ACTION MATRIX
% ============================================================

\section*{Action Matrix}

\begin{tabularx}{\linewidth}{@{}p{0.10\linewidth}p{0.10\linewidth}p{0.20\linewidth}Y p{0.10\linewidth}@{}}
\toprule
\textbf{ID} & \textbf{Severity} & \textbf{File} & \textbf{Required change} & \textbf{Effort} \\
\midrule
REF-001 & major & paper/main.tex & Harmonize rare-class language throughout with abstract hedging. & small \\
REF-002 & major & paper/main.tex & Narrow novelty claim; acknowledge Wu et al.\ per-class results. & small \\
REF-003 & major & paper/main.tex & Add Cuoco et al.\ 2025; update Wu et al.\ to CQG. & small \\
REF-004 & major & paper/main.tex & Apply consistent speculative hedging to all interpretations. & small \\
REF-005 & major & paper/main.tex & Add DC difference to abstract; downgrade ``first quantitative.'' & small \\
REF-006 & minor & paper/main.tex & Correct 0.670 to 0.669. & trivial \\
REF-007 & minor & paper/main.tex & Fix ``8 rare classes'' label. & trivial \\
REF-008 & minor & paper/main.tex & Soften section heading. & trivial \\
REF-009 & minor & paper/main.tex & Motivate 20\% threshold or remove. & trivial \\
REF-010 & minor & paper/main.tex & Add imbalanced-learning citations. & trivial \\
REF-011 & minor & paper/main.tex & Add ViT interpretability citations. & trivial \\
REF-012 & minor & table\_overall.tex & Switch to two-sided test or add footnote. & trivial \\
REF-013 & minor & paper/main.tex & Update Wu et al.\ to CQG published version. & trivial \\
\bottomrule
\end{tabularx}

% ============================================================
% CONFIDENCE MATRIX
% ============================================================

\section*{Confidence Matrix}

\begin{tabularx}{\linewidth}{@{}p{0.26\linewidth}p{0.16\linewidth}Y@{}}
\toprule
\textbf{Dimension} & \textbf{Confidence} & \textbf{Notes} \\
\midrule
Correctness & HIGH & All key quantities independently recomputed from source JSON. \\
Novelty & HIGH & Literature search verified; Wu et al.\ comparison confirmed. \\
Reproducibility & HIGH & Manifest independently reviewed. \\
Publishability & MEDIUM & Final CQG editorial judgment depends on referee pool composition. \\
\bottomrule
\end{tabularx}

\vfill

{\color{ReportCrimson}\rule{\linewidth}{1.1pt}}

\small
\textit{This is an AI-generated mock referee report. It supplements but does not replace expert peer review.}

\end{document}
